% Section content template for individual Educative lessons
% This file contains the actual content for a specific section
% Generated from Educative API section content

% Section: Quiz on Twitter's Design
% Section ID: 6740284133605376
% Section Slug: quiz-on-twitters-design
% Generated: 2025-12-08T12:07:53.069952

\subsection*{Quiz}
\vspace{0.3cm}
\noindent
\textbf{Question 1:} Suppose you work as an engineer at a well-known company providing social services. Your task is to improve real-time event processing to get frequent results. Which tool below will help you the most in this scenario?
\\[0.2cm]
\begin{itemize}
 \item BigQuery
\\[0.1cm]
\textit{Explanation:} 
BigQuery is a serverless, fully managed data warehouse that allows for scalable analysis of vast amounts of data.
 \item \textbf{[CORRECT]} Apache Kafka
\\[0.1cm]
\textit{Explanation:} 
Apache Kafka is a framework for distributed real-time event processing and storage.
 \item Vertica
\\[0.1cm]
\textit{Explanation:} 
Vertica is a columnar data storage that can manage massive amounts of data and provides lightning-fast query performance in usually time-consuming circumstances.
 \item FlockDB
\\[0.1cm]
\textit{Explanation:} 
FlockDB is a distributed, fault-tolerant graph database for handling large network graphs.
\end{itemize}
\vspace{0.3cm}
\noindent
\textbf{Question 2:} Assume that an account of a public figure with millions of followers posts a Tweet on Twitter. Your system will need to handle millions of interactions with this Tweet by the account's millions of followers. Which solution below is the best way to handle this traffic, assuming that the system knows the audience details of the Tweet?
\\[0.2cm]
\begin{itemize}
 \item Dedicate a few servers to the public figure in the central data center(s).
\\[0.1cm]
\textit{Explanation:} 
Increase burden on central data center and also increase latency.
 \item \textbf{[CORRECT]} Increase or dedicate some servers to particular locations where a large portion of the audience belongs.
 \item Increase the number of servers in all Twitter's data centers situated in different regions to handle massive traffic.
\\[0.1cm]
\textit{Explanation:} 
If most of the traffic is coming from a few regions, increasing servers in other datacenters is useless and resource wastage.
 \item Show the Tweet to followers region-wise to avoid handling massive traffic.
\\[0.1cm]
\textit{Explanation:} 
All users having connections with tweets must get the tweet. So, it is not fulfilling the service.
\end{itemize}
\vspace{0.3cm}
\noindent
\textbf{Question 3:} When the user sends a home timeline request, the system responds with a specified number of Top-k tweets to the user. On what basis does the system identify the Top-k tweets for a particular user?
\\[0.2cm]
\begin{itemize}
 \item The user liked the Tweets (posted by the unfollowed account) and Trends.
\\[0.1cm]
\textit{Explanation:} 
Trends metric is irrelevant here.
 \item The user's followers and location.
\\[0.1cm]
\textit{Explanation:} 
Users will not get the tweets posted by their followers.
 \item \textbf{[CORRECT]} The user's followed accounts' Tweets and Retweeted Tweets by the followed accounts.
 \item Location and unfollowed accounts
\\[0.1cm]
\textit{Explanation:} 
Users will not get unfollowed accounts' tweets as long as tweets are not promoted.
\end{itemize}
\vspace{0.3cm}
\noindent
\textbf{Question 4:} Sharded counters are an efficient solution to count multiple interactions (like or view) on thousands of celebrities' Tweets individually. The problem is that when the user interacts with any Tweet, it needs an immediate response (for instance, an increase or decrease in the like count). Which solution below will help in this regard?
\\[0.2cm]
\begin{itemize}
 \item Maintain counters in the server cache.
\\[0.1cm]
\textit{Explanation:} 
It is difficult to manage millions of counters in the cache because they are limited in size (memory).
 \item \textbf{[CORRECT]} Place specified shards of the particular counters near the end users based on the celebrity's history (Last few Tweets' audiences).
\\[0.1cm]
\textit{Explanation:} 
In this case, the nearest server will immediately respond to the end-users who liked or viewed the tweet.
 \item Increase the number of shards of multiple counters that belong to the Tweet in the central data center.
\\[0.1cm]
\textit{Explanation:} 
It will increase latency because central data centers are distant from the end-users and also increase the central servers' burden.
 \item Create a real-time counter in the client cache with the specified detail of the Tweet.
\\[0.1cm]
\textit{Explanation:} 
This solution has multiple problems, such as limited client cache. For instance, if the end-user performs many operations (like, view, reply) on many tweets, creating and managing many counters in the cache is difficult. Second , what if the client liked a post and the relevant counter is incremented in the client cache. But somehow client cache gets empty (due to any failure) before moving that details to the application server. So
, this solution is not reliable.
\end{itemize}
\vspace{0.3cm}
\noindent
\textbf{Question 5:} \textbf{(Select all that apply.)} We discussed how the client-side load balancers work and what kind of benefit it provides. Which of the options below were identified as possible drawbacks of client-side load balancers?
\\
\textit{(Select all that apply.)}
\\[0.2cm]
\begin{itemize}
 \item They act as a single point of failure.
 \item \textbf{[CORRECT]} They're not suitable for small-scale applications because there are no dedicated servers for the services.
\\[0.1cm]
\textit{Explanation:} 
Each service is managed within the backend server. Therefore, no such type of communication (server selection service to service) is required among the server.
 \item They increase network hop.
 \item \textbf{[CORRECT]} They make it more difficult to update all clients when servers add/drop or algorithm modify.
\end{itemize}



% End of section content